\section{The project folder and its configuration}
\label{sec:folder}

Nothing needs moving or renaming: every path in the code is relative to the \texttt{final\_thesis}
folder. \texttt{notebooks\textbackslash} holds the eight pipeline notebooks,
\texttt{config.yaml} and \texttt{requirements.txt}. \texttt{data\textbackslash raw\textbackslash}
holds the collected review histories and a checkpoint per game;
\texttt{data\textbackslash processed\textbackslash} holds everything derived from them, and
\texttt{data\textbackslash ml\_runs\textbackslash} the MLflow database.
\texttt{web\textbackslash} holds the Astro site, the Streamlit application and the script that
builds the data bundle both read. \texttt{figures\textbackslash} and
\texttt{reports\textbackslash} hold the charts and the two documents.


Every path, threshold, band cutoff and chart colour lives in
\texttt{notebooks\textbackslash config.yaml}, which the notebooks, the website and the demo
application all read. A value written anywhere else would be a defect, and that is why the report,
the site and the model cannot disagree about a cutoff. Table~\ref{tab:config} covers the blocks
worth knowing before you run anything.

\begin{table}[htbp]
\centering
\caption{The settings that matter in \texttt{config.yaml}.}
\label{tab:config}
\small
\begin{tabular}{|p{2.6cm}|p{11.4cm}|}
\hline
\textbf{Block} & \textbf{What it controls} \\
\hline
\texttt{seed} & fixed at 42, used by every random step so a re-run reproduces the same numbers \\
\hline
\texttt{collection\_date} & 2026-08-01, the day playtime was snapshotted. Changing it re-collects
nothing \\
\hline
\texttt{paths} & where each stage reads and writes \\
\hline
\texttt{frame} & which games are modelled: at least 200 English early-access reviews and 50 after
launch, giving 545 \\
\hline
\texttt{split} & temporal, on launch date, 70 per cent training \\
\hline
\texttt{bands} & the At Risk and Healthy cutoffs, 43.6 and 57.2, written back by stage 7 from the
training data. Never type these by hand \\
\hline
\texttt{predictors} & the explicit list of columns the model may use \\
\hline
\texttt{palette} & every colour used in the charts, the site and the application \\
\hline
\end{tabular}
\end{table}

\shot{config_yaml}{0.85}{\texttt{config.yaml}, showing the paths, the modelling frame and the
calibrated band cutoffs. The notebooks run in the order given in
Section~\ref{sec:pipeline}, not the alphabetical order the folder shows.}

\textbf{One thing not to change.} \texttt{data\textbackslash raw\textbackslash} is written once and
never edited. Every correction happens downstream of it, so any result traces back to the bytes
that came off the endpoint.
